\chapter{問題点と解決提案手法}
\label{sec:approach}

\section{評価手順に残る問題}
評価資料には，判定が十分に付与されていない文書や，特定の分野に偏った検索要求が含まれる場合がある．
そのような資料を用いた比較では，得点が何を表しているかを慎重に解釈する必要がある．
手法の有効性と，資料の構成によって生じる差を完全に切り分けることは容易ではない．
まず，確認できた制約を記録し，結論の範囲を明確にすることが重要である．

実装上の問題も比較の信頼性に影響する．入力の欠落，要求識別子の重複，文書の正規化処理の違いなどがあると，
意図した条件で比較できなくなる．これらの問題は，結果の平均値だけを見ても発見しにくい．
評価前に入力と出力の対応を確認する処理を設けることで，不完全な結果をそのまま集計することを防げる．

\section{提案する確認手順}
本例で提案するのは，検索要求ごとの対応関係を保存し，集計前に欠落と重複を点検する手順である．
両方の結果がそろっていることを確認した後に差を求め，群ごとの集計へ進む．
欠落がある場合には，その扱いを評価規則として明示する．自動的に得点をゼロに置き換える処理は，
その規則が意図したものか確認した上で採用する必要がある．

この手順では，比較条件と評価結果を別々のファイルとして保存し，実行識別子で結び付ける．
条件は何を行う予定だったかを示し，実行ログは実際に何が起きたかを示す．
予想外の挙動があった場合にも，その情報を捨てずに残すことが，再現と原因の調査に役立つ．
ここで述べた手順は説明用の提案であり，効果を実証した結果ではない．

% Demonstration-only break.
\newpage
\section{検証方法}
提案した確認手順自体も点検する必要がある．少数の検索要求を用いた例で，
入力識別子の対応，得点の計算，平均値への集計を手作業で追えるようにする．
その上で，欠落や重複を含む入力に対して意図した動作になるかを確認する．
正常な入力で動くだけでは，誤った入力を検出できることまでは保証されない．

処理の再現性を確認する際には，乱数の種だけでなく，ソフトウェアと資料の版も記録する．
並列処理や外部サービスに依存する部分がある場合には，完全に同じ結果が得られない要因を説明する．
再実行で変化する範囲が分かれば，その変化と手法による差を区別して検討しやすくなる．

\section{今後の課題}
次の段階として，異なる特徴を持つ検索要求の集合で，同じ比較を実施することが考えられる．
追加する資料は，単に規模を増やすためではなく，現在の解釈を制約している不確実性を調べるために選ぶ．
元の実験と共通にする条件と，変更が必要な条件を区別すれば，得られた違いを説明しやすくなる．

計算時間やメモリ使用量も検討対象となり得る．順位の改善が得られたとしても，
応答時間の増加が利用場面に適しているとは限らない．測定に含める処理の範囲を定め，
実行環境とともに結果を報告する必要がある．どの課題から取り組むかは，
研究の主張を支える上で最も重要な未解決点が何かに基づいて判断する．
